% =====================================================================
% 07_portfolio_engine.tex — Appendix B.2: pricing kernels, portfolios,
% risk premia (second-order portfolio engine) (WP-D)
% =====================================================================
\section{The portfolio and risk-premium engine (second-order)}
\label{sec:portfolio-engine}

The first-order engine of Section~\ref{sec:engine} determines prices, returns,
and the realized wealth-share revaluation \emph{given} the steady-state
portfolio $(\qk k,b_F,b^g_{H,s})$. It does not pin that portfolio down: to first
order, all assets are perfect substitutes and the portfolio is indeterminate.
KL close the model with the Devereux--Sutherland (2011) device: the
\emph{equilibrium} portfolio is the one that makes the \emph{second-order}
optimal-portfolio (excess-return) conditions hold, and to leading order it is
determined by a first-order risk-sharing condition. This section reconstructs
KL's online Appendix~B.2 (their app.~pp.~73--75). The outputs are: (i) the
second-order portfolio-Euler conditions; (ii) the log pricing-kernel deviation
$\dev{\pk}_{0,1}$ and its Home and Foreign ``risk-adjusted'' forms
$\dev{\pk}_{0,1}-\rra\tel^z\dev{\epsilon}^z_1$ and
$\fr{\dev{\pk}}_{0,1}-\rraF\tel^z\dev{\epsilon}^z_1$; (iii) the
local-completeness (efficient-risk-sharing) condition; and (iv) the
international-portfolio equation, KL's eq.~(99), with its composite coefficient
$\Gam$. These feed directly into Propositions~\ref{prop:portfolios}
and~\ref{prop:riskpremium}.

Throughout we follow KL's app.~B.2 in writing $\tel^z$ for the loading of the
productivity surprise (their $\sigma^z$; macro \texttt{\textbackslash tel}$^z$ in
the engine), $\dev{\epsilon}^z_1$ for the period-$1$ productivity innovation, and
$\lwedge$ for the steady-state \emph{labor wedge}. As in the engine we keep
$\hbias$ general until completeness is imposed.

\subsection{The Devereux--Sutherland portfolio-Euler conditions}
\label{sec:portfolio-euler}

A second-order approximation of the period-$0$ optimal portfolio-choice
conditions (their app.~p.~73) gives, for the excess return of capital over the
safe dollar bond and for the excess return of the Foreign bond over the safe
dollar bond,
\begin{align}
\Ezero\bigl[\rk_1-\rr_1\bigr] + \text{Jensen terms}
&= \dev{\cy}_0 - \Ezero\Bigl[\bigl(\dev{\pk}_{0,1}-\rra\tel^z\dev{\epsilon}^z_1\bigr)\bigl(\rk_1-\rr_1\bigr)\Bigr], \label{eq:ds-capital}\\
\Ezero\bigl[\fr{r}_1-\D\dev{\rer}_1-\rr_1\bigr] + \text{Jensen terms}
&= \dev{\cy}_0 - \Ezero\Bigl[\bigl(\dev{\pk}_{0,1}-\rra\tel^z\dev{\epsilon}^z_1\bigr)\bigl(\fr{r}_1-\D\dev{\rer}_1-\rr_1\bigr)\Bigr], \label{eq:ds-fbond}
\end{align}
and analogously for the Foreign household (with $\fr{\dev{\pk}}_{0,1}$ and
$\rraF$ in place of $\dev{\pk}_{0,1}$ and $\rra$). Here $\dev{\pk}_{0,1}$ is the
log deviation of the Home household's real pricing kernel (period~$0$ to
period~$1$), defined in Subsection~\ref{sec:log-kernel} below, and
$\dev{\cy}_0$ is the period-$0$ (steady-state) convenience yield on the safe
dollar bond, which appears because the dollar bond yields liquidity services that
capital and Foreign bonds do not.

\paragraph{Reading the conditions.} The left-hand side of each equation is an
expected excess return plus a \emph{Jensen} (variance-of-returns) correction. KL
note (their app.~p.~73) that ``\emph{Jensen terms} reflect the component of
excess returns which do not reflect safety shocks nor the covariance with agents'
pricing kernels (and instead reflect the variance of returns).'' The
economically active part of each condition is therefore the right-hand side: the
convenience premium $\dev{\cy}_0$ that the safe bond commands, net of the
covariance of the (risk-adjusted) pricing kernel with the excess return. The
equilibrium portfolio is the one that makes these covariances take the values
required for \eqref{eq:ds-capital}--\eqref{eq:ds-fbond} to hold. Because the
Jensen terms are pure second moments of returns that are common across the
household problems and do not load on the safety shock or on the kernel
covariance, KL's strategy is to read the equilibrium holdings off the
\emph{first-order} risk-sharing condition that the covariance terms encode; the
Jensen terms then determine only the (common, second-order) level of expected
excess returns, not the cross-sectional portfolio. We make this precise via the
local-completeness condition of Subsection~\ref{sec:local-completeness}.

\subsection{The log pricing kernel}
\label{sec:log-kernel}

In period~$1$, the log deviation of the representative Home household's pricing
kernel is (their app.~p.~73)
\begin{equation}
\dev{\pk}_{0,1} = -\dev{\rs c}_1 + (1-\rra)\,\dev{\rs \val}_1, \label{eq:m01}
\end{equation}
which is the log-linearization of the Epstein--Zin/recursive stochastic discount
factor: with unit IES (the simplified environment), the kernel is
$m_{0,1}\propto (c_1)^{-1}(v_1/ce_0)^{1-\rra}$, whose log deviation is
$-\dev{\rs c}_1$ plus $(1-\rra)$ times the (re-scaled) value deviation
$\dev{\rs\val}_1$ (the certainty-equivalent normalizer $ce_0$ is non-stochastic
as of period~$0$ and drops from the period-$0$-to-$1$ innovation).
\gapflag{KL display \eqref{eq:m01} directly as ``the log deviation in the
representative Home household's pricing kernel is given by
$\dev{\pk}_{0,1}=-\dev{\rs c}_1+(1-\rra)\dev{\rs\val}_1$.'' We supply the
one-line provenance: under unit IES the recursive kernel is the product of the
consumption-growth term (exponent $-1$ on $c$, since $1/\text{IES}=1$) and the
continuation-value/certainty-equivalent term carrying the risk-aversion exponent
$1-\rra$. The macro renders $\dev{\rs c}_1=\hat{\tilde c}_1$ and
$\dev{\rs\val}_1=\hat{\tilde v}_1$, the period-$1$ re-scaled (stationarized)
consumption and value deviations of Section~\ref{sec:setup}.}

\paragraph{The value linearization.} The (re-scaled) recursive value linearizes
to (their app.~p.~73)
\begin{equation}
\dev{\rs \val}_1 = (1-\disc)\dev{\rs c}_1 - (1-\disc)(1-\lwedge)(1-\cshare)\dev{\ell}_1 + \disc\,\dev{\rs{ce}}_1, \label{eq:v1}
\end{equation}
where $\lwedge$ is the labor wedge in the deterministic steady state and
$\dev{\rs{ce}}_1$ is the log deviation of the certainty equivalent of the
period-$2$ continuation value. The term
$-(1-\disc)(1-\lwedge)(1-\cshare)\dev{\ell}_1$ is the flow disutility of labor,
valued at the steady-state \emph{after-wedge} marginal rate; the factor
$(1-\lwedge)$ scales the labor term down by the wedge because, with a positive
labor wedge, the private marginal disutility of work is below the social marginal
product, so a unit of employment contributes
$(1-\lwedge)(1-\cshare)$ to flow utility (in consumption units) rather than the
full $(1-\cshare)$.
\gapflag{KL display \eqref{eq:v1} as
$\dev{\rs\val}_1=(1-\disc)\dev{\rs c}_1-(1-\disc)(1-\lwedge)(1-\cshare)\dev{\ell}_1+\disc\dev{\rs{ce}}_1$.
The labor coefficient $(1-\lwedge)(1-\cshare)$ is KL's; we supply the reading.
Note this is the only place the steady-state labor wedge $\lwedge$ enters the
kernel, and it is precisely the channel that makes \emph{employment} (not only
productivity and wealth) a priced state variable in Props~4--5; it vanishes if
$\lwedge=0$.}

\paragraph{The certainty equivalent.} Because there is no uncertainty after
period~$1$, the period-$1$ certainty equivalent equals the period-$2$ expected
value (their app.~p.~74),
\begin{equation}
\dev{\rs{ce}}_1 = \Eone\dev{\rs \val}_2 = \atil\,\Eone\dev{\wsh}_2 + \cshare\,\widetilde{\widetilde{\dev k}}_1, \label{eq:ce1}
\end{equation}
where the second equality is the engine's period-$2$ value
relation~\eqref{eq:Ev2} (their app.~eq.~(77)). To first order the
certainty-equivalent operator is the identity here, since the only remaining
randomness (the period-$1$ innovations) is already realized at the point where
$\dev{\rs{ce}}_1$ is evaluated.

\paragraph{Assembling $\dev{\rs c}_1$ from the engine.} The period-$1$ Home
consumption is the engine's reduced form~\eqref{eq:c1-solved}, which we restate
here in the app.~B.2 notation (writing the recurring denominator from the engine
explicitly rather than as $\Lambda$, to match KL's displayed lines):
\begin{equation}
\dev{\rs c}_1 = \atil\,\Eone\dev{\wsh}_2
+(1-\cshare)\Bigl(\tfrac{1}{1+\zs}\dev{\ell}_1+\tfrac{\zs}{1+\zs}\fr{\dev{\ell}}_1\Bigr)
+\cshare\,\widetilde{\widetilde{\dev k}}_1
-\frac{\hbias}{\frac{\cshare}{1-\cshare}+\tel+(1-\tel)\left(\hbias\frac{1+\zs}{\zs}\right)^2}\bigl(\fr{\dev{\ell}}_1-\dev{\ell}_1\bigr). \label{eq:c1-b2}
\end{equation}
Equation~\eqref{eq:c1-b2} is identically \eqref{eq:c1-solved} with the
denominator $\Lambda$ \eqref{eq:Lambda-def} written out; KL display it in this
expanded form on app.~p.~73.

\paragraph{The combined value deviation.} Substituting \eqref{eq:c1-b2} and
\eqref{eq:ce1} into the value linearization~\eqref{eq:v1} and collecting terms,
KL obtain (their app.~p.~74)
\begin{multline}
\dev{\rs \val}_1 = \atil\,\Eone\dev{\wsh}_2 - \cshare\,\tel^z\dev{\epsilon}^z_1
+\lwedge(1-\disc)(1-\cshare)\dev{\ell}_1\\
+(1-\disc)\left((1-\cshare)\frac{\zs}{1+\zs}
-\frac{\hbias}{\frac{\cshare}{1-\cshare}+\tel+(1-\tel)\left(\hbias\frac{1+\zs}{\zs}\right)^2}\right)\bigl(\fr{\dev{\ell}}_1-\dev{\ell}_1\bigr). \label{eq:v1-combined}
\end{multline}
\gapflag{\eqref{eq:v1-combined} is KL's displayed combined value line (app.~p.~74).
The reconstruction: write \eqref{eq:v1} as
$\dev{\rs\val}_1=(1-\disc)\dev{\rs c}_1-(1-\disc)(1-\lwedge)(1-\cshare)\dev{\ell}_1+\disc\dev{\rs{ce}}_1$.
The $\disc\dev{\rs{ce}}_1=\disc[\atil\Eone\dev{\wsh}_2+\cshare\widetilde{\widetilde{\dev k}}_1]$
piece and the $(1-\disc)\atil\Eone\dev{\wsh}_2$ piece inside $(1-\disc)\dev{\rs c}_1$
combine to the full $\atil\Eone\dev{\wsh}_2$ (coefficients $\disc+(1-\disc)=1$);
likewise the two $\cshare\widetilde{\widetilde{\dev k}}_1$ pieces combine to
$\cshare\widetilde{\widetilde{\dev k}}_1$. KL then \emph{re-express} the symmetric
employment piece using the impact wage-setting relation
\eqref{eq:predet-wage}: under wages set one period in advance the common labor
movement satisfies $(1-\cshare)(\frac{1}{1+\zs}\dev{\ell}_1+\frac{\zs}{1+\zs}\fr{\dev{\ell}}_1)+\cshare\widetilde{\widetilde{\dev k}}_1$
contributes a $-\cshare\tel^z\dev{\epsilon}^z_1$ term once the productivity
surprise is substituted (the same $\tel^z\dev{\epsilon}^z_1$ that loads the wage
equation), so the aggregate-employment + capital block collapses into the
$-\cshare\tel^z\dev{\epsilon}^z_1$ displayed term. The labor coefficient
collects as follows: $(1-\disc)\dev{\rs c}_1$ contributes
$+(1-\disc)(1-\cshare)\frac{1}{1+\zs}$ to $\dev{\ell}_1$ and the relative-labor
$-\hbias/[\cdots]$ term; the explicit disutility term contributes
$-(1-\disc)(1-\lwedge)(1-\cshare)\dev{\ell}_1$. Adding the
$+(1-\disc)(1-\cshare)\dev{\ell}_1$ implicit in the symmetric piece, the
$\dev{\ell}_1$ coefficient nets to $+\lwedge(1-\disc)(1-\cshare)$ (the
$-(1-\lwedge)$ and $+1$ combine to $+\lwedge$), which is the displayed
$\lwedge(1-\disc)(1-\cshare)\dev{\ell}_1$ wedge term; the residual relative-labor
piece collects into the displayed
$(1-\disc)\bigl((1-\cshare)\frac{\zs}{1+\zs}-\hbias/[\cdots]\bigr)(\fr{\dev{\ell}}_1-\dev{\ell}_1)$
coefficient. This collection is the single tersest algebra step in B.2; we have
reproduced KL's displayed coefficients and flag that we filled the term-by-term
collection rather than re-deriving each cancellation from the primitive Euler
block.}

\paragraph{The Home risk-adjusted kernel.} Substituting \eqref{eq:c1-b2} and
\eqref{eq:v1-combined} into the kernel definition~\eqref{eq:m01}, and grouping the
productivity loading on the kernel as $\rra\tel^z\dev{\epsilon}^z_1$, KL obtain
the central B.2 expression (their app.~p.~74)
\begin{multline}
\dev{\pk}_{0,1} - \rra\,\tel^z\dev{\epsilon}^z_1
= -\rra\Bigl[\atil\,\Eone\dev{\wsh}_2 + (1-\cshare)\tel^z\dev{\epsilon}^z_1\Bigr]
- (1-\cshare)\Bigl(\tfrac{1}{1+\zs}\dev{\ell}_1+\tfrac{\zs}{1+\zs}\fr{\dev{\ell}}_1\Bigr)\\
+\frac{\hbias}{\frac{\cshare}{1-\cshare}+\tel+(1-\tel)\left(\hbias\frac{1+\zs}{\zs}\right)^2}\bigl(\fr{\dev{\ell}}_1-\dev{\ell}_1\bigr)\\
+(1-\rra)(1-\disc)\left[\lwedge(1-\cshare)\dev{\ell}_1
+\left((1-\cshare)\frac{\zs}{1+\zs}-\frac{\hbias}{\frac{\cshare}{1-\cshare}+\tel+(1-\tel)\left(\hbias\frac{1+\zs}{\zs}\right)^2}\right)\bigl(\fr{\dev{\ell}}_1-\dev{\ell}_1\bigr)\right]. \label{eq:m01-home}
\end{multline}
\gapflag{\eqref{eq:m01-home} is KL's displayed Home risk-adjusted-kernel line
(app.~p.~74). The reconstruction is purely mechanical given \eqref{eq:m01},
\eqref{eq:c1-b2}, \eqref{eq:v1-combined}: $\dev{\pk}_{0,1}=-\dev{\rs c}_1+(1-\rra)\dev{\rs\val}_1$.
The first bracketed group $-\rra[\atil\Eone\dev{\wsh}_2+(1-\cshare)\tel^z\dev{\epsilon}^z_1]$
arises because the $\atil\Eone\dev{\wsh}_2$ term appears in $\dev{\rs c}_1$ with
coefficient $-1$ and in $(1-\rra)\dev{\rs\val}_1$ with coefficient $(1-\rra)$,
netting to $-\rra\atil\Eone\dev{\wsh}_2$; we then \emph{move} the
$\rra\tel^z\dev{\epsilon}^z_1$ loading to the left-hand side (KL's convention,
isolating the part of the kernel orthogonal to the directly-priced productivity
risk), which converts the $-\cshare\tel^z\dev{\epsilon}^z_1$ pieces of
$\dev{\rs c}_1$ and $\dev{\rs\val}_1$ into the displayed
$-\rra(1-\cshare)\tel^z\dev{\epsilon}^z_1$ (using $-(-\cshare)-(1-\rra)\cshare+\rra=\rra(1-\cshare)$
on the $\tel^z\dev{\epsilon}^z_1$ coefficient). The symmetric-employment and
relative-labor terms from $-\dev{\rs c}_1$ appear with the signs shown, and the
$(1-\rra)(1-\disc)[\cdots]$ block is exactly $(1-\rra)$ times the
wedge-and-relative-labor part of \eqref{eq:v1-combined}. We reproduce KL's line
term-for-term; the only judgement call is the left-hand-side grouping of
$\rra\tel^z\dev{\epsilon}^z_1$, which is KL's own (it is the quantity that enters
the Devereux--Sutherland conditions \eqref{eq:ds-capital}--\eqref{eq:ds-fbond}).}

\paragraph{The Foreign risk-adjusted kernel.} Analogous steps in Foreign
(replacing $\rra\to\rraF$, the wealth-share loading $\atil\to-\frac{1}{\zs}\atil$
from the period-$2$ Foreign value~\eqref{eq:Ev2-star}, and the symmetric-labor
weight as appropriate) yield (their app.~p.~74)
\begin{multline}
\fr{\dev{\pk}}_{0,1} - \rraF\,\tel^z\dev{\epsilon}^z_1
= -\rraF\Bigl[-\tfrac{1}{\zs}\atil\,\Eone\dev{\wsh}_2 + (1-\cshare)\tel^z\dev{\epsilon}^z_1\Bigr]
- (1-\cshare)\Bigl(\tfrac{1}{1+\zs}\dev{\ell}_1+\tfrac{\zs}{1+\zs}\fr{\dev{\ell}}_1\Bigr)\\
-\frac{1}{\zs}\frac{\hbias}{\frac{\cshare}{1-\cshare}+\tel+(1-\tel)\left(\hbias\frac{1+\zs}{\zs}\right)^2}\bigl(\fr{\dev{\ell}}_1-\dev{\ell}_1\bigr)\\
+(1-\rraF)(1-\disc)\left[\lwedge(1-\cshare)\fr{\dev{\ell}}_1
-\left((1-\cshare)\frac{1}{1+\zs}-\frac{1}{\zs}\frac{\hbias}{\frac{\cshare}{1-\cshare}+\tel+(1-\tel)\left(\hbias\frac{1+\zs}{\zs}\right)^2}\right)\bigl(\fr{\dev{\ell}}_1-\dev{\ell}_1\bigr)\right]. \label{eq:m01-foreign}
\end{multline}
\gapflag{\eqref{eq:m01-foreign} is KL's displayed Foreign risk-adjusted-kernel
line (app.~p.~74). The reconstruction mirrors the Home case under the Foreign
relabeling: the Foreign household's period-$2$ value carries
$-\frac{1}{\zs}\atil\Eone\dev{\wsh}_2$ (\eqref{eq:Ev2-star}), so its
wealth-share loading flips sign and scales by $1/\zs$; the relative-labor
coefficient picks up the $1/\zs$ that distinguishes the Foreign cross-country
weight; the wedge term carries $\fr{\dev{\ell}}_1$ rather than $\dev{\ell}_1$.
We reproduce KL's displayed signs; the structural mirror to \eqref{eq:m01-home}
is exact, and we flag that we transcribed the Foreign weights ($1/\zs$ factors,
$\frac{1}{1+\zs}$ vs.\ $\frac{\zs}{1+\zs}$) from KL's line rather than
re-deriving the Foreign value linearization from scratch.}

\subsection{Local completeness and the risk-sharing condition}
\label{sec:local-completeness}

KL observe (their app.~p.~74) that the simplified environment is \emph{locally
complete} in the sense of Coeurdacier and Gourinchas (2016): with three traded
assets (safe dollar bond, capital, Foreign bond) and two shocks (safety
$\dev{\cy}$, productivity $\dev{\epsilon}^z$), the asset span is rich enough that
the equilibrium portfolios implement constrained-efficient risk sharing to first
order. The implication is the \emph{risk-sharing condition}
\begin{equation}
\boxed{\;\dev{\pk}_{0,1} - \rra\,\tel^z\dev{\epsilon}^z_1
= \fr{\dev{\pk}}_{0,1} - \rraF\,\tel^z\dev{\epsilon}^z_1 + \D\dev{\rer}_1.\;} \label{eq:risk-sharing}
\end{equation}
This is the open-economy Backus--Smith / law-of-one-price relation for stochastic
discount factors: the Home and Foreign risk-adjusted kernels must agree up to the
real-exchange-rate change $\D\dev{\rer}_1$, which converts Foreign marginal
utility into Home units. Equivalently, the two countries' marginal valuations of
state-contingent wealth coincide once expressed in a common numeraire --- the
hallmark of (locally) complete markets.

\paragraph{Why this pins the portfolio.} The kernels \eqref{eq:m01-home} and
\eqref{eq:m01-foreign} are functions of the endogenous outcomes
$(\Eone\dev{\wsh}_2,\dev{\ell}_1,\fr{\dev{\ell}}_1)$, which in turn depend on the
portfolio through the realized wealth-share identity \eqref{eq:wsh-realized}
(their eq.~(92)) and the labor block. Imposing \eqref{eq:risk-sharing} therefore
delivers one scalar equation linking the portfolio weights to the shocks; this is
the equation that the international-portfolio relation below makes explicit. The
Jensen terms of \eqref{eq:ds-capital}--\eqref{eq:ds-fbond} do not enter
\eqref{eq:risk-sharing} because they are common across the Home and Foreign
conditions and cancel in the difference --- which is precisely why KL can read the
equilibrium portfolio off the first-order risk-sharing condition rather than the
full second-order excess-return level.
\gapflag{KL state \eqref{eq:risk-sharing} as the local-completeness implication
``$\dev{\pk}_{0,1}-\rra\tel^z\dev{\epsilon}^z_1=\fr{\dev{\pk}}_{0,1}-\rraF\tel^z\dev{\epsilon}^z_1+\D\dev{\rer}_1$''
citing Coeurdacier and Gourinchas (2016) for the local-completeness notion. We
supply the interpretation (Backus--Smith with $\D\dev{\rer}_1$ as the numeraire
conversion) and the explanation of why imposing it pins the portfolio while the
Jensen terms drop. We do not re-derive the Coeurdacier--Gourinchas completeness
result; we take it as KL do, namely as the statement that three assets span the
two shocks plus the redistributive (wealth-share) margin.}

\subsection{The international-portfolio equation (their eq.~99)}
\label{sec:eq99}

Imposing the risk-sharing condition~\eqref{eq:risk-sharing} --- i.e.\ subtracting
\eqref{eq:m01-foreign} from \eqref{eq:m01-home} and setting the difference equal
to $\D\dev{\rer}_1$ --- and substituting the realized wealth-share
identity~\eqref{eq:wsh-realized} for $\Eone\dev{\wsh}_2$ via $\dev{\wsh}_1$
(their app.~eq.~(92)), KL collect terms to obtain the
\emph{international-portfolio equation} (their app.~eq.~(99)):
\begin{multline}
\Bigl(\frac{\qk k}{\sav}-1\Bigr)\bigl(\rk_1-\rr_1\bigr)
+\frac{b_F}{\sav}\bigl(\fr{r}_1-\dev{\rer}_1-\rr_1\bigr)
-\disc\,\frac{\zs}{1+\zs}\,\frac{b^g_{H,s}}{\sav}\,\dev{\cy}_1
=\\
\frac{1}{\Gam}(\rraF-\rra)(1-\cshare)\tel^z\dev{\epsilon}^z_1
+\frac{1}{\Gam}(\rraF-\rra)\lwedge(1-\disc)(1-\cshare)\Bigl(\tfrac{1}{1+\zs}\dev{\ell}_1+\tfrac{\zs}{1+\zs}\fr{\dev{\ell}}_1\Bigr)\\
+(1-\disc)\frac{\zs}{1+\zs}(\tel-1)\frac{1-\cshare}{\cshare}
\frac{1-\left(\hbias\frac{1+\zs}{\zs}\right)^2}{\frac{\cshare}{1-\cshare}+\tel+(1-\tel)\left(\hbias\frac{1+\zs}{\zs}\right)^2}\bigl(\fr{\dev{\ell}}_1-\dev{\ell}_1\bigr)\\
-\frac{1}{\Gam}\Bigl(\tfrac{1}{\zs}(\rraF-1)+(\rra-1)\Bigr)(1-\disc)(1-\lwedge)(1-\cshare)\frac{\zs}{1+\zs}\bigl(\fr{\dev{\ell}}_1-\dev{\ell}_1\bigr)\\
+\frac{1}{\Gam}\Bigl(\tfrac{1}{\zs}(\rraF-1)+(\rra-1)\Bigr)(1-\disc)
\frac{\hbias}{\frac{\cshare}{1-\cshare}+\tel+(1-\tel)\left(\hbias\frac{1+\zs}{\zs}\right)^2}\bigl(\fr{\dev{\ell}}_1-\dev{\ell}_1\bigr), \label{eq:eq99}
\end{multline}
where the composite coefficient is
\begin{equation}
\Gam \;\equiv\; \left[\rra + \frac{1}{\zs}\rraF
+\frac{\hbias^2\left(\frac{1+\zs}{\zs}\right)^3}{\frac{\cshare}{1-\cshare}+\tel\left(1-\left(\hbias\frac{1+\zs}{\zs}\right)^2\right)}\right]\atil. \label{eq:Gamma-def}
\end{equation}
\gapflag{Equations~\eqref{eq:eq99}--\eqref{eq:Gamma-def} are KL's displayed
eq.~(99) and the definition of $\Gam$ (their app.~p.~75). This is the single
densest derivation in B.2 and KL display only the final collected form
(``Substituting in using the above results and those of the previous section and
collecting terms, we obtain''). We reconstruct the structure as follows. (i) The
\emph{left-hand side} is exactly the realized wealth-share
identity~\eqref{eq:wsh-realized}/(92): the three portfolio-weighted excess-return
terms (leverage, Foreign bond, seigniorage) equal $\dev{\wsh}_1=\Eone\dev{\wsh}_2$
on impact (recall \eqref{eq:Ewsh-solved} sets
$\Eone\dev{\wsh}_2=\dev{\wsh}_1+(\text{relative-labor})$). KL move all of
\eqref{eq:wsh-realized} to the LHS so the RHS is the risk-sharing
\emph{requirement} on $\Eone\dev{\wsh}_2$. (ii) The \emph{right-hand side} is the
difference of the two risk-adjusted kernels~\eqref{eq:m01-home}$-$\eqref{eq:m01-foreign}
minus $\D\dev{\rer}_1$, solved for $\Eone\dev{\wsh}_2$. The
$\Eone\dev{\wsh}_2$ coefficient on that difference is
$-\rra\atil-(-\rraF(-\frac{1}{\zs}\atil))=-(\rra+\frac{1}{\zs}\rraF)\atil$
\emph{plus} the $\D\dev{\rer}_1$ contribution: $\D\dev{\rer}_1=\dev{\rer}_1$
carries an $\Eone\dev{\wsh}_2$ loading through \eqref{eq:q1-solved}, namely
$\hbias\frac{1+\zs}{\zs}\cdot\frac{\hbias(\frac{1+\zs}{\zs})^2}{\frac{\cshare}{1-\cshare}+\tel(1-(\hbias\frac{1+\zs}{\zs})^2)}\atil
=\frac{\hbias^2(\frac{1+\zs}{\zs})^3}{\frac{\cshare}{1-\cshare}+\tel(1-(\hbias\frac{1+\zs}{\zs})^2)}\atil$,
which is the third term inside the bracket of $\Gam$. Collecting,
$\Gam=[\rra+\frac{1}{\zs}\rraF+\frac{\hbias^2(\frac{1+\zs}{\zs})^3}{\frac{\cshare}{1-\cshare}+\tel(1-(\hbias\frac{1+\zs}{\zs})^2)}]\atil$
is exactly the coefficient multiplying $\Eone\dev{\wsh}_2$ in the
kernel-difference; dividing through by it gives the $1/\Gam$ on the RHS terms.
(iii) The RHS \emph{numerators} are the remaining (non-$\Eone\dev{\wsh}_2$) terms
of \eqref{eq:m01-home}$-$\eqref{eq:m01-foreign}: the
$(\rraF-\rra)(1-\cshare)\tel^z\dev{\epsilon}^z_1$ productivity term (from the
$-\rra(1-\cshare)$ vs.\ $-\rraF(1-\cshare)$ loadings on $\tel^z\dev{\epsilon}^z_1$);
the $(\rraF-\rra)\lwedge(1-\disc)(1-\cshare)(\cdots)$ aggregate-employment term
(from the wedge terms, which carry $(1-\rra)$ vs.\ $(1-\rraF)$ and net to
$(\rraF-\rra)\lwedge$ on the symmetric-employment piece); and the
relative-labor $(\fr{\dev{\ell}}_1-\dev{\ell}_1)$ terms, which split into the
$(\tel-1)$ ``terms-of-trade revaluation'' piece (matching the $\Eone\dev{\wsh}_2$
relative-labor coefficient of \eqref{eq:Ewsh-solved}, carried over because the
LHS used \eqref{eq:wsh-realized} while the RHS prices $\Eone\dev{\wsh}_2$) and the
$(\frac{1}{\zs}(\rraF-1)+(\rra-1))$ pieces from the $(1-\rra)$/$(1-\rraF)$
relative-labor coefficients. We have matched KL's displayed eq.~(99)
term-by-term: the five RHS lines correspond, in order, to the productivity term,
the aggregate-employment term, the $(\tel-1)$ terms-of-trade-revaluation term,
and the two relative-labor wedge/home-bias terms. We flag this as a \emph{genuine
reconstruction}: KL display only the collected result, and while every
coefficient below matches their eq.~(99), we filled the intermediate
kernel-difference algebra and the identification of the $\Gam$ bracket with the
$\Eone\dev{\wsh}_2$ loading on $\dev{\rer}_1$. The Props~4--5 proofs use only the
\emph{structure} of \eqref{eq:eq99} (which RHS terms vary with $\dev{\epsilon}^z$
vs.\ $\dev{\cy}$, and the signs of their coefficients), not the exact algebra of
every term.}

\paragraph{The three drivers of Home wealth.} The left-hand side of
\eqref{eq:eq99} is the realized revaluation of Home's wealth share
(Proposition~\ref{prop:wealth}); the right-hand side is what efficient risk
sharing \emph{requires} that revaluation to be. KL read off three economic
drivers (their app.~p.~75); international risk sharing calls for Home wealth (the
LHS) to rise with:
\begin{itemize}[leftmargin=1.6em,itemsep=3pt]
  \item \textbf{productivity, provided $\rraF>\rra$}: a positive TFP shock
  $\dev{\epsilon}^z_1>0$ raises aggregate production and hence world consumption;
  the more risk-tolerant country (lower risk aversion) should bear more of this
  aggregate risk, so if Home is more risk-tolerant ($\rra<\rraF$) Home wealth
  should rise in good productivity states. This is the
  $\frac{1}{\Gam}(\rraF-\rra)(1-\cshare)\tel^z\dev{\epsilon}^z_1$ term, positive
  when $\rraF>\rra$.
  \item \textbf{aggregate employment, provided $\lwedge(\rraF-\rra)>0$}: with a
  positive labor wedge, an increase in employment raises welfare (the marginal
  product exceeds the private marginal disutility), and again the more
  risk-tolerant country should be exposed to this welfare gain. This is the
  $\frac{1}{\Gam}(\rraF-\rra)\lwedge(1-\disc)(1-\cshare)(\cdots)$ term.
  \item \textbf{Foreign employment less Home employment}, i.e.\ the relative-labor
  gap $\fr{\dev{\ell}}_1-\dev{\ell}_1$, under either of two stated sign
  conditions (their app.~p.~75): (a)
  $(\tel-1)\bigl(1-(\hbias\frac{1+\zs}{\zs})^2\bigr)>0$, ``since this implies that
  Foreign labor income rises relative to Home labor income''; or (b)
  $-\bigl(\frac{1}{\zs}(\rraF-1)+(\rra-1)\bigr)\bigl(\frac{\hbias/[\cdots]}{1}-(1-\lwedge)(1-\cshare)\frac{\zs}{1+\zs}\bigr)>0$,
  ``since this implies that Home requires more wealth when its real exchange rate
  appreciates, despite the larger disutility of labor in Foreign.''
\end{itemize}
KL close B.2 with the statement that the equilibrium risk premium on Foreign
bonds relative to dollar bonds is ``the covariance of (the negative of) the log
deviation in any agent's pricing kernel with the excess log return,'' i.e.\
\begin{equation}
\text{risk premium on Foreign vs.\ dollar bonds}
\;=\; -\Cov_0\bigl(\dev{\pk}_{0,1}-\rra\tel^z\dev{\epsilon}^z_1,\;\fr{r}_1-\D\dev{\rer}_1-\rr_1\bigr), \label{eq:rp-cov}
\end{equation}
which is the object Proposition~\ref{prop:riskpremium} signs. (By local
completeness~\eqref{eq:risk-sharing} the same covariance obtains using the
Foreign household's kernel, so the risk premium is well-defined independent of
whose kernel one prices with.)
